% Generated draft: requires mathematical and source review.
\documentclass[11pt]{article}
\usepackage[margin=25mm]{geometry}
\usepackage{fontspec}
\setmainfont{Latin Modern Roman}
\usepackage{amsmath,amssymb,mathtools,unicode-math}
\setmathfont{Latin Modern Math}
\usepackage{enumitem,needspace,xurl,hyperref}
\setlength{\parindent}{0pt}
\setlength{\parskip}{5pt}
\setlength{\emergencystretch}{3em}
\newcommand{\N}{\mathbb{N}}
\newcommand{\Z}{\mathbb{Z}}
\newcommand{\Q}{\mathbb{Q}}
\newcommand{\R}{\mathbb{R}}
\newcommand{\C}{\mathbb{C}}
\newcommand{\F}{\mathbb{F}}
\newcommand{\A}{\mathbb{A}}
\newcommand{\PP}{\mathbb{P}}
\newcommand{\E}{\mathbb{E}}
\newcommand{\eps}{\varepsilon}
\newcommand{\dd}{\,\mathrm{d}}
\begin{document}
\section*{UnsolvedMath ID: 701725}
\textbf{Source problem code:} catalogue entry
\subsection{Definitions and mathematical statement}
Let \(G\) be a finite group, let \(p\) be a prime, and let \(B\) be a Brauer \(p\)-block of \(G\), with defect group \(D\). Write \(\operatorname{Irr}(B)\) for the set of ordinary irreducible complex characters belonging to \(B\), and define
\[
\operatorname{cd}(B)=\{\chi(1):\chi\in\operatorname{Irr}(B)\}.
\]
For a finite group \(H\), write \(\operatorname{Irr}(H)\) for its irreducible complex characters and set
\[
\operatorname{cd}(H)=\{\chi(1):\chi\in\operatorname{Irr}(H)\}.
\]
These are sets of distinct degrees, without multiplicity. Define the derived series of \(D\) by \(D^{(0)}=D\) and \(D^{(i+1)}=[D^{(i)},D^{(i)}]\) for integers \(i\geq 0\). Since \(D\) is a finite \(p\)-group, its derived length is
\[
\operatorname{dl}(D)=\min\{n\in\Z_{\geq 0}:D^{(n)}=1\}.
\]
In particular, the trivial group has derived length zero.

\begin{enumerate}
\item[(a)] Is it true that, for every such \(G,p,B,D\),
\[
\operatorname{dl}(D)\leq |\operatorname{cd}(B)|?
\]
Here the source phrase ``bounded by'' is read as ``at most.''
\item[(b)] (A. Jaikin-Zapirain.) Is it true that, for every such \(G,p,B,D\),
\[
|\operatorname{cd}(D)|\leq |\operatorname{cd}(B)|?
\]
\end{enumerate}
\subsection{Short English statement}
For every Brauer block of a finite group, does the number of distinct ordinary irreducible character degrees in the block bound the derived length of its defect group? Part (b), attributed to A. Jaikin-Zapirain, asks whether it also bounds the number of distinct irreducible complex character degrees of the defect group.
\subsection{Sources}
\begin{itemize}
\item Local UnsolvedMath registry, permanent ID \(701725\), exact source code \texttt{catalogue entry}, title \emph{Kourovka Notebook Problem 20.67}. The supplied record labels the problem open; this label is supplied metadata, not an independently verified assessment. The record was created and updated on 2026-10-16.
\item \emph{The Kourovka Notebook (Unsolved Problems in Group Theory)}, problem 20.67, proposed by G. Navarro in 2022. The supplied background identifies the 20th issue (2022) and the 21st issue edition (2026), with identifier \texttt{arXiv:1401.0300v48}. The supplied collection description names E. I. Khukhro and V. D. Mazurov as editors.
\item Part (b) is specifically attributed in the supplied statement to A. Jaikin-Zapirain.
\end{itemize}
\end{document}
